\section{F0: Does the Reward Bit Change Written Memory?}
\label{sec:f0}
We begin with a paired intervention over released WebArena Shopping trajectories from the Salesforce fragility study. A deterministic sampler requests six source trajectories, balanced across original benchmark outcomes. For each trajectory, the task text and full action history are copied byte-for-byte into two ReasoningBank writer calls. One call uses the released success-reflection instruction and the paired call uses the released failure-reflection instruction. DeepSeek-V4-Flash writes memory at temperature zero, and raw provider responses are content-addressed before analysis.

Four of the six requested trajectory pairs return complete memory outputs in both conditions. The two incomplete requests terminate in the failure-label writer before assistant text is produced and are excluded from the paired effect calculation. Among the four complete pairs, every pair changes normalized memory content and the extracted memory-item title set. Token-set Jaccard distance is 0.691 for task 21 and 0.708 for task 22. It is 0.770 for both task 23 and task 25, giving mean distance 0.735. The 4/4 change rate is therefore an existence result conditional on paired completion. Panel (a) of Figure~\ref{fig:write-control} shows the complete pair-level result.

The content changes are procedural as well as lexical. For task 21, success-conditioned memory centers on entering the review section and preserving direct quotations, whereas failure-conditioned memory emphasizes exhaustive pagination and extraction verification. For task 23, the success writer emphasizes evidence collection across pages while the failure writer emphasizes pagination state and broader synonym coverage. To make this distinction explicit without introducing another learned judge, we reuse a deterministic operation-slot taxonomy over the frozen memory-item titles: navigation, pagination/coverage, query expansion, semantic matching, evidence capture, extraction, verification/completeness, temporal filtering, and status filtering. The slot set changes in all four complete pairs; mean slot-set Jaccard distance is 0.493 (range 0.400--0.571). Moreover, every complete pair has a success-only slot in evidence capture, navigation, or query expansion, while verification/completeness appears as a failure-only slot in three of four pairs. This is a structural diagnostic rather than embedding-level semantic proof; Section~\ref{sec:prompt-control} asks whether the same structure also changes under reward-preserving prompt rewording.

The F0 result establishes the first link of the mechanism: changing the reward-conditioned writer branch changes persistent memory under a fixed experience. Because the branch itself changes the writer instruction, the next section tests the strongest simple alternative explanation: generic sensitivity to prompt wording.
